\chapter{線形代数I：空間・写像・行列}
\label{ch:tb-linear-one}

\section{部分空間、基底、次元}

\begin{definition}{部分空間}{tb-subspace}
実または複素ベクトル空間 $V$ の空でない部分集合 $W$ が\textbf{部分空間}であるとは
\[
u,v\in W,\ \alpha,\beta\in\mathbb F
\Longrightarrow \alpha u+\beta v\in W
\]
を満たすことをいう。$0\in W$ と加法・スカラー倍についての閉性を別々に確認してもよい。
\end{definition}

\begin{proposition}{共通部分と合併}{tb-subspace-intersection-union}
部分空間 $U,W\subset V$ に対し $U\cap W$ は部分空間である。
$U\cup W$ が部分空間となるための必要十分条件は
$U\subset W$ または $W\subset U$ である。
\end{proposition}

\begin{proof}
共通部分では線形結合が $U,W$ の両方に残る。
合併について、どちらの包含も成り立たないなら
$u\in U\setminus W$, $w\in W\setminus U$ を取れる。
$u+w\in U\cup W$ と仮定すると、$u+w\in U$ なら $w\in U$、
$u+w\in W$ なら $u\in W$ となり、いずれも矛盾する。
\end{proof}

\begin{definition}{生成・一次独立・基底}{tb-basis}
$v_1,\ldots,v_m$ の一次結合全体を
\[
\operatorname{span}\{v_1,\ldots,v_m\}
=\left\{\sum_{i=1}^m a_iv_i\ \middle|\ a_i\in\mathbb F\right\}
\]
と書く。$\sum a_iv_i=0\Rightarrow a_1=\cdots=a_m=0$ のとき一次独立という。
一次独立で $V$ を生成する組を基底といい、有限次元空間の基底の本数を $\dim V$ とする。
\end{definition}

基底であることは「一次独立」と「生成」の二項目を示す。
行列の列ベクトルなら、pivot 列が像の基底、自由変数の本数が核の次元を与える。

\begin{definition}{和と直和}{tb-direct-sum}
\[
U+W=\{u+w\mid u\in U,w\in W\}.
\]
$U\cap W=\{0\}$ のとき $U+W$ を $U\oplus W$ と書く。
この条件は表示 $v=u+w$ の一意性と同値である。
\end{definition}

\begin{theorem}{次元公式}{tb-dimension-formula}
有限次元部分空間 $U,W$ に対し
\[
\dim(U+W)=\dim U+\dim W-\dim(U\cap W).
\]
特に直和なら $\dim(U\oplus W)=\dim U+\dim W$ である。
\end{theorem}

\section{線形写像、核、像}

\begin{definition}{線形写像と核・像}{tb-linear-map}
$f:U\to V$ が
$f(\alpha x+\beta y)=\alpha f(x)+\beta f(y)$ を満たすとき線形写像という。
\[
\ker f=\{x\in U\mid f(x)=0\},\qquad
\operatorname{Im}f=\{f(x)\mid x\in U\}.
\]
どちらも部分空間である。
\end{definition}

\begin{theorem}{単射判定と次元定理}{tb-rank-nullity}
\[
f\text{ が単射}\Longleftrightarrow\ker f=\{0\}.
\]
$U$ が有限次元なら
\[
\dim U=\dim\ker f+\dim\operatorname{Im}f.
\]
同じ有限次元空間の間の線形写像では、単射、全射、全単射が同値である。
\end{theorem}

行列 $A$ が表す写像では、$A$ の pivot 列が $\operatorname{Im}A$ の基底を与え、
$Ax=0$ の一般解から $\ker A$ の基底を得る。
二つの像の共通部分は
\[
Au=Bv\Longleftrightarrow
\begin{pmatrix}A&-B\end{pmatrix}
\binom uv=0
\]
を解き、得られた $Au$ から一次独立なものを選べばよい。

\begin{proposition}{同型写像は基底を基底へ送る}{tb-isomorphism-basis}
$f:V\to V$ が線形同型で $v_1,\ldots,v_n$ が基底なら
$f(v_1),\ldots,f(v_n)$ も基底である。
\end{proposition}

\begin{proof}
$\sum a_if(v_i)=0$ なら線形性により $f(\sum a_iv_i)=0$ である。
単射性から $\sum a_iv_i=0$、基底の一次独立性から全ての $a_i=0$ となる。
本数が $\dim V$ に等しいため、これは基底である。
\end{proof}

\begin{proposition}{像が直和である二写像}{tb-two-map-direct-sum}
$f,g:\mathbb F^n\to\mathbb F^n$ が
$\mathbb F^n=\operatorname{Im}f\oplus\operatorname{Im}g$ を満たすなら
\[
\ker(f+g)=\ker f\cap\ker g.
\]
さらに $\mathbb F^n=\ker f+\ker g$ なら $f+g$ は全単射である。
\end{proposition}

\begin{proof}
$(f+g)x=0$ なら $f(x)=-g(x)$ は二つの像の共通部分に属する。
直和なので両方0となる。逆包含は直ちに従う。
さらに次元定理と像の直和から
\[
\dim\ker f+\dim\ker g
=2n-\dim\operatorname{Im}f-\dim\operatorname{Im}g=n.
\]
一方、核の和が全空間なら次元公式により
\[
n=\dim(\ker f+\ker g)
=n-\dim(\ker f\cap\ker g).
\]
従って $\ker(f+g)=\ker f\cap\ker g=\{0\}$ であり、$f+g$ は全単射である。
\end{proof}

\section{行基本変形と連立方程式}

行基本変形は
\begin{enumerate}
  \item 二行を交換する、
  \item 一行を0でない定数倍する、
  \item 一行に別の行の定数倍を加える
\end{enumerate}
の三つであり、連立方程式の解集合を変えない。

拡大係数行列を階段形にし、
\[
\operatorname{rank}A=\operatorname{rank}(A\mid b)
\]
なら $Ax=b$ は解を持つ。pivot のない列に自由変数を置き、全ての解を書く。
斉次方程式 $Ax=0$ の解空間は $\ker A$ であり
\[
\dim\ker A=n-\operatorname{rank}A.
\]
正則行列の逆行列は Gauss--Jordan 法
$(A\mid I)\longrightarrow(I\mid A^{-1})$ で同時に求められる。

$A$ が正則なら、行列方程式は掛ける側に注意して
\[
XA=B\Longrightarrow X=BA^{-1},\qquad
AY=C\Longrightarrow Y=A^{-1}C
\]
と解く。

\section{三角行列と二段階解法}

\begin{proposition}{前進代入・後退代入}{tb-triangular-systems}
対角成分が全て0でない下三角行列 $L$ に対する $Ly=b$ は、
第1式から順に $y_1,y_2,\ldots$ を決める\textbf{前進代入}で一意に解ける。
上三角行列 $U$ に対する $Ux=y$ は、最終式から逆順に決める\textbf{後退代入}で解ける。
\end{proposition}

$A=LU$ と分解できれば
\[
Ax=b\Longleftrightarrow L(Ux)=b
\]
なので、まず $Ly=b$、次に $Ux=y$ を解けばよい。また
\[
\det A=\det L\det U
=\left(\prod_i l_{ii}\right)\left(\prod_i u_{ii}\right).
\]
特に $L$ の対角成分が全て1なら $\det A$ は $U$ の対角成分の積である。

$L$ と $A$ から $A=LU$ の $U$ を求めるときも前進代入と同じ順序を使う。
$L$ の対角成分が1なら、行ベクトルを $A_i,U_i$ と書いて
\[
U_i=A_i-\sum_{k<i}l_{ik}U_k
\]
と上の行から順に決定できる。

\section{行列式と正則性}

\begin{proposition}{行列式の計算規則}{tb-determinant-properties}
\[
\det(AB)=\det A\det B,\quad
\det A^{\mathsf T}=\det A,\quad
\det(cA)=c^n\det A.
\]
行交換で符号が反転し、一行の $c$ 倍で行列式も $c$ 倍、
一行に他の行の倍を加えても行列式は変わらない。
三角行列の行列式は対角成分の積である。
\end{proposition}

\begin{theorem}{正則性の同値条件}{tb-invertibility}
$n$ 次正方行列 $A$ について次は同値である。
\[
A\text{ は正則},\quad \det A\ne0,\quad \operatorname{rank}A=n,
\quad\ker A=\{0\},\quad Ax=b\text{ は任意の }b\text{ に一意解を持つ}.
\]
\end{theorem}

行列式は存在・非存在の証明にも使う。実行列で $A^2=-I_n$ なら
\[
(\det A)^2=\det(-I_n)=(-1)^n.
\]
左辺は非負なので $n$ は偶数である。$n=2$ では
$\begin{pmatrix}0&-1\\1&0\end{pmatrix}$ が例になる。
一方、任意の実ベクトル $x$ に対し
$x^{\mathsf T}A^{\mathsf T}Ax=\norm{Ax}^2\ge0$ なので
$A^{\mathsf T}A=-I_n$ は不可能である。

\section{行列多項式から逆行列を作る}

\begin{proposition}{逆行列を式から読む}{tb-polynomial-inverse}
$A^2+A-2I=O$ なら
\[
A(A+I)=2I,\qquad A^{-1}=\frac12(A+I).
\]
一般に $Aq(A)=cI$（$c\ne0$）まで変形できれば $A^{-1}=q(A)/c$ である。
また $N^m=O$ なら
\[
(I-N)^{-1}=I+N+\cdots+N^{m-1},\qquad
(I+N)^{-1}=I-N+\cdots+(-N)^{m-1}.
\]
\end{proposition}

積を実際に掛けて $I-N^m=I$ または $I-(-N)^m=I$ になることを示せば、
正則性と逆行列を同時に証明できる。
